\documentstyle[%


righttag%


,11pt%


,amssymb%


,russian%        remove in Eng.


,izvru%         Set izven% in Eng.


]{amsart}





% {\leq} 624


\newcommand{\Mod}{\operatorname{Mod}}


\newcommand{\Hom}{\operatorname{Hom}}


\newcommand{\End}{\operatorname{End}}


\newcommand{\Alg}{\operatorname{Alg}}


\newcommand{\Map}{\operatorname{Map}}


\newcommand{\e}{\varepsilon}


\newcommand{\Fr}{\operatorname{Fr}}


\newcommand{\J}{{\frak G}}


\newcommand{\R}{{\frak R}}


\newcommand{\K}{{\frak K}}


\newcommand{\M}{{\frak M}}


\newcommand{\N}{{\frak N}}


\newcommand{\U}{{\frak U}}


\newcommand{\E}{{\frak E}}


\newcommand{\F}{{\frak F}}


\newcommand{\r}{{\frak r}}


\renewcommand{\O}{{\frak O}}


\newcommand{\ep}{\varepsilon}


\newcommand{\Fb}{{\frak F}}


\newcommand{\sym}{\Sigma}


\newcommand{\s}{\sigma}


\renewcommand{\t}{\tau}


\newcommand{\al}{\alpha}


\newcommand{\be}{\beta}


\newcommand{\la}{\lambda}


\renewcommand{\o}{\omega}


\newcommand{\fb}{{\bold f}}


\newcommand{\tn}{\otimes}


\newcommand{\T}{{\frak T}}


\newcommand{\f}{\varphi}


\newcommand{\Fi}{\Phi}


\newcommand{\p}{\psi}


\newcommand{\Ps}{\Psi}


\newcommand{\A}{{\frak A}}


\newcommand{\G}{{\Gamma}}


\newcommand{\ar}[1]{\overset{#1}{\longrightarrow}{}}


\newcommand{\old}[1]{\ov{\ov{#1}}}


\newcommand{\ti}[1]{\widetilde{#1}}


\newcommand{\tts}[1]{}








%\hoffset=-15mm%         For Eng.


\setcounter{page}{50}


\god{2004}


\nomer{4~(503)} %            Remove in Eng.


%\nomer{Vol.\,48, No.\,4}%        Set in Eng.


%\str{\pageref{begin}--\pageref{end}}% Set in Eng.


\udk{519.175:512.565}





\author{\it С.Н.\,Тронин, А.В.\,Семенова}


% NB:  ^^ remove in Eng.


\title{Операды конечных помеченных графов}





\date{}


%\pagestyle{myheadings}%     Set in Eng.


\pagestyle{plain} %       Remove in Eng.


\begin{document}


%\thispagestyle{plain}%     Set in Eng.


\maketitle


\label{begin}








Данная работа посвящена изучению некоторых структур операд,


возникающих на множествах помеченных конечных графов.


Определения и обозначения, относящиеся к теории операд,


взяты из \cite{TK}. Структуры,


изучаемые в данной работе, были определены в \cite{T}. Для конечных


непомеченных графов композиция графов, очень похожая на операдную,


введена в \cite{LM}. Для непомеченных графов, однако,


нельзя использовать возможности, предоставляемые теорией операд.


Используемые понятия и обозначения из теории графов соответствуют книге


\cite{H}. Заметим, что термин ``операда графов'' недавно уже появился


в литературе по теории операд \cite{Oz}, но  в данной статье


этот термин используется для совсем иного объекта.


Ввиду того что элементы наших операд


больше похожи на графы в традиционном смысле, сохраним слово ``графы'' в


названии изучаемого здесь объекта. Авторы благодарны проф. З.\,Озиевичу


за предоставленную возможность ознакомиться с его работами,


в частности, со статьей \cite{Oz}.








Пусть $K$ --- некоторое множество с выделенным элементом


$\varepsilon$. Обозначим через $\M_{n,k}=\M_{n,k}(K)$ множество всех


$n\times k$-матриц c компонентами из $K$. Определим две операции композиции.


Пусть $A_i\in \M_{n_i,k_i}$, $1\leq i\leq m$, $B\in \M_{m,m}$,


$B=(b_{ij})$, $1\leq j\leq m$.


Операция композиции, которую будем называть {\it композицией}


1, определяется так:


\begin{equation}


A_1 A_2 \dotsc A_m B=


\begin{pmatrix}


A_1 & b_{1\,2} & \dotsc & b_{1\,m} \\


b_{2\,1} & A_2 & \dotsc & b_{2\, m} \\


\vdots & \vdots  & \ddots & \vdots \\


b_{m\,1} & b_{m\,2} & \dotsc & A_m 


\end{pmatrix}


,


\end{equation}


$A_1 A_2 \dotsc A_m B$ --- блочная


$(n_1+\dotsb +n_m)\times (k_1+\dotsb +k_m)$-матрица, 


разбитая на блоки размерами $n_i\times k_j$, и


$b_{ij}$ в (1) --- матрица размера $n_i\times k_j$,


заполненная одним и тем же элементом $b_{ij}$ из $B$.


Вместо $b_{ij}$ следовало бы писать $b_{ij}J_{n_i, k_j}$, где


$J_{n_i,k_j}$ --- матрица, состоящая из единиц.





Пусть теперь $K$ --- моноид с единицей $\varepsilon$.


Операция композиции, которую будем называть {\it композицией} 2, 


определяется так:


\begin{equation}


A_1 A_2 \dotsc A_m B=


\begin{pmatrix}


A_1b_{1\,1} & b_{1\,2} & \dotsc & b_{1\,m} \\


b_{2\,1} & A_2b_{2\,2} & \dotsc & b_{2\, m} \\


\vdots & \vdots  & \ddots & \vdots \\


b_{m\,1} & b_{m\,2} & \dotsc & A_mb_{m\,m} \\


\end{pmatrix}


.


\end{equation}


Соглашения здесь те же самые, что и в (1).





Напомним \cite{TK}, что операция композиции называется ассоциативной,


если выполнены тождества


$$


(\ov{A}_1 B_1)\dotsc (\ov{A}_m B_m)C=


(\ov{A}_1\dotsc \ov{A}_m)(B_1\dotsc B_mC),


$$


где


$\overline{A}_i=A_{1,i}\dotsc A_{n_i,i}$, $A_{l,i}\in \M_{n_{l,i},k_{l,i}}$,


$B_i\in \M_{n_i,n_i}$, $C\in \M_{m,m}$,


$1\leq i\leq m$, $ 1\leq l\leq n_i$.








\begin{lem}


Обе операции композиции $(1)$ и $(2)$ ассоциативны.


\end{lem}





\begin{pf}


Это делается с помощью непосредственной проверки,


использующей определения (1) и (2).


\end{pf}








Определим два семейства множеств


$\M=\{ \M(n) \mid n=1,2, \dotsc \}$,


$\ov{\M}=\{ \ov{\M}(n) \mid n=1,2, \dotsc \}$. Пусть


$\M(n)$ --- подмножество $\M_{n,n}$, состоящее из всех


матриц, диагональные элементы которых равны $\varepsilon$, а


$\ov{\M}(n)=\M_{n,n}$.


На $\M(n)$ и $\ov{\M}(n)$ слева действует группа подстановок $n$-й степени


$\sym_n$ следующим образом.


Если $A=(a_{ij})$, $1\leq i,j\leq n$, $a_{ij}\in K$, $\s\in \sym_n$,


то матрица $\s A$ состоит из элементов


\begin{equation}


(\s A)_{ij}=a_{\s^{-1}(i),\s^{-1}(j)}.


\end{equation}


В дальнейшем при действиях с матрицами можно считать, что их компоненты лежат в 


полугрупповом кольце ${\bold Z}[K]$.








\begin{lem}


Пусть $M(\s )$ --- матрица подстановки $\s \in \sym_n$. Тогда


$$


\s A=M(\s )A{M(\s )}^{-1}.


$$


\end{lem}





\begin{pf}


Пусть $E_{i,j}$ --- матричная единица,


тогда $M(\s)=\sum\limits_{j=1}^n E_{\s (j),j}$, и


\begin{gather*}


M(\s )AM(\s ^{-1})=\bigg(\sum\limits_{k=1}^n E_{\s (k),k}\bigg)


\bigg(\sum\limits_{i,j=1}^n a_{ij}E_{i,j}\bigg)


\bigg(\sum\limits_{l=1}^n E_{l,\s (l)}\bigg)=


\bigg(\sum\limits_{j=1}^n \sum\limits_{k=1}^n a_{kj}E_{\s (k),j}\bigg)


\bigg(\sum\limits_{l=1}^n E_{l,\s (l)}\bigg)=\\


%


=\sum\limits_{k=1}^n \sum\limits_{l=1}^n a_{kl}E_{\s (k),\s (l)}=


\sum\limits_{i=1}^n \sum\limits_{j=1}^n a_{\s ^{-1}(i),\s ^{-1}(j)}E_{i,j}.


\end{gather*}


Получена матрица, в которой на $(i,j)$-м месте стоит элемент


$a_{\s ^{-1}(i),\s ^{-1}(j)}$, как и у $\s A$.


\end{pf}





\begin{thm}


$1)$ Семейство $\M$ $($ соответственно $\ov{\M})$ с операцией композиции $1$


$($соот\-вет\-ст\-вен\-но~$2)$ и действием симметрических групп $(3)$


становится операдой.





$2)$


Множество $\A=\bigcup \limits_{n,k\geq 1} \M_{n,k}$ превращается в алгебру 


над операдой $\M$, если определить


композицию по формуле $(1)$, и в алгебру над операдой $\ov{\frak M}$, если 


определить композицию по формуле $(2)$.


При этом операды $\M$ и $\ov{\M}$ надо считать несимметрическими 


$($т.\,е. исключить из их определения 


действия симметрических групп$)$.


\end{thm}





\begin{pf}


Ассоциативность проверена в лемме 1. Единичным элементом в $\M$ и $\ov{\M}$ 


cлужит $1\times 1$-матрица


$(\varepsilon)\in \M (1)=\ov {\M} (1)$, где


$\varepsilon\in K$ --- выделенный элемент $K$ (или единица моноида).


Соответствующие свойства сразу вытекают


из формул (1), (2).


Все это показывает, что $\M$ и $\ov {\M}$ --- несимметрические операды.


Из леммы 1 теперь следует, что


$\A$ есть алгебра над $\M$ или $\ov {\M}$ в зависимости от того,


какая композиция используется.





Остается проверить, что формула (3) превращает $\M$ и $\ov {\M}$ в


симметрические операды. Доказательства


для $\M$ и $\ov{\M}$ практически одинаковы. Выполним все проверки для $\M$.


Докажем тождество


\begin{equation}


(\s_1A_1)(\s_2A_2)\dotsc (\s_mA_m)B=


(\s_1\times\s_2\times\dotsb\times \s_m)A_1A_2\dotsc A_mB.


\end{equation}


Будем использовать лемму 2. Напомним, что 


$$


\s\times\t=


\begin{pmatrix}


1 & \dotsc & m & m+1 & \dotsc & m+n \\


\s (1) & \dotsc & \s (m) & m+\t (1) & \dotsc & m+\t (n)


\end{pmatrix}


,


$$


где $\s\in\sym_m$, $\t\in\sym_n$. Тогда


$M(\s\times\t)=\left(


\begin{smallmatrix}M(\s ) & 0 \\


0 & M(\t )\end{smallmatrix}\right)$.


Легко проверяется равенство


\begin{gather*}


\begin{pmatrix}


M(\s _1)A_1M(\s _1)^{-1} & b_{12} & \dotsc & b_{1m}\\


b_{21} & M(\s _2)A_2M(\s _2)^{-1} & \dotsc & b_{2m}\\


\dotsc & \dotsc & \ddots & \dotsc\\


b_{m1} & b_{m2} & \dotsc & M(\s _m)A_mM(\s _m)^{-1}


\end{pmatrix}


= \\


%


=


\begin{pmatrix}


M(\s _1) & \dotsc & 0\\


0 & \dotsc & 0\\


\dotsc & \ddots & \dotsc\\


0 & \dotsc & M(\s _m)


\end{pmatrix}


\begin{pmatrix}


A_1 & b_{12}J_{12} & \dotsc & b_{1m}J_{1m}\\


b_{21}J_{21} & A_2 & \dotsc & b_{2m}J_{2m}\\


\dotsc & \dotsc & \ddots & \dotsc\\


b_{m1}J_{m1} & b_{m2}J_{m2} & \dotsc & A_m


\end{pmatrix}


\begin{pmatrix}


M(\s _1)^{-1} & \dotsc & 0\\


0 & \dotsc & 0\\


\dotsc & \ddots & \dotsc\\


0 & \dotsc & M(\s _m)^{-1}


\end{pmatrix}


,


\end{gather*}


из которого следует (4).





Осталось проверить выполнимость последнего тождества


\begin{equation}


A_1A_2\dotsc A_m(\s B)=


(\al * \s)(A_{\s (1)}A_{\s (2)}\dotsc A_{\s (m)}B).


\end{equation}


Напомним \cite{TK}, что $\al * \s$ определяется так: 


если $\al=(n_1,\dotsc,n_m)$, $\s\in \sym_m$,


$$


{\bold p}_1=(1, 2, \dotsc ,n_1), {\bold p}_2=(n_1+1, \dotsc ,n_1+n_2), \dotsc ,


{\bold p}_i=(n_1+\dotsb +n_{i-1}+1,\dotsc ,n_1+\dotsb +n_{i-1}+n_i),


$$


$1\leq i\leq m$, то


$\al * \s =


\begin{pmatrix}


{\bold p}_1 & {\bold p}_2 & \dotsc & {\bold p}_m\\


{\bold p}_{\s (1)} & {\bold p}_{\s (2)} & \dotsc & {\bold p}_{\s (m)}


\end{pmatrix}$. Очевидно, 


$$


A_1\dotsc A_m(\s B)=


\begin{pmatrix}


A_1 & b_{\s ^{-1}(1),\s ^{-1}(2)} & \dotsc & b_{\s ^{-1}(1),\s ^{-1}(m)}\\


b_{\s ^{-1}(2),\s ^{-1}(1)} & A_2 & \dotsc & b_{\s ^{-1}(2),\s ^{-1}(m)}\\


\dotsc & \dotsc & \ddots & \dotsc\\


b_{\s ^{-1}(m),\s ^{-1}(1)} & b_{\s ^{-1}(m),\s ^{-1}(2)} & \dotsc & A_m


\end{pmatrix}


.


$$





Рассмотрим правую часть тождества (5).


Матрица $M(\al * \s)$ устроена следующим образом.


Это блочная матрица из $m \times m$ блоков,


где блок с индексами $(\s (j),j)$ имеет $n_{\s (j)}$ строк и $n_{\s (j)}$


столбцов и равен


единичной $n_{\s (j)}\times n_{\s (j)}$-матрице $E_{n_{\s (j)}}$.


Остальные блоки, размеры которых


однозначно определяются, нулевые. 


Неформально говоря, $M(\al * \s )$ получается так: в $M(\s )$ 


вместо единиц (расположенных в строках $\s (j)$


и столбцах $j$ ) ``подставляются'' блоки $E_{n_{\s (j)}}$, а вместо


нулей --- нулевые блоки соответствующих размеров.


Пусть


$$


A_{\s (1)}\dotsc A_{\s (m)}B=


\begin{pmatrix}


A_{\s (1)} & b_{12} & \dotsc & b_{1m}\\


b_{21} & A_{\s (2)} & \dotsc & b_{2m}\\


\dotsc & \dotsc & \ddots & \dotsc \\


b_{m1} & b_{m2} & \dotsc & A_{\s (m)}


\end{pmatrix}


=C=


\begin{pmatrix}


C_{11} & C_{12} & \dotsc & C_{1m}\\


C_{21} & C_{22} & \dotsc & C_{2m}\\


\dotsc & \dotsc & \ddots & \dotsc \\


C_{m1} & C_{m2} & \dotsc & C_{mm}


\end{pmatrix}


, 


$$


где $C_{ij}$ --- блок размера $n_{\s (i)}\times n_{\s (j)}$.


Вычислим $M(\al * \s)CM(\al * \s)^{-1}=


(\al * \s)C$ с учетом того, что $M(\al * \s)^{-1}= M(\al * \s)^T$,


т.\,е. в $j$-й блочной строке шириной $n_{\s (j)}$,


в $\s (j)$-м блочном столбце ширины $n_{\s (j)}$ 


находится матрица $E_{n_{\s (j)}}$. Положим


$$


M=M(\al * \s )=


\begin{pmatrix}


M_{11} & \dotsc & M_{1m} \\


\dotsc & \ddots & \dotsc \\


M_{m1} & \dotsc & M_{mm} 


\end{pmatrix}


,


$$


где $M_{ij}$ --- блок размером $n_i\times n_{\s (j)}$, причем


$M_{\s (j),j}=E_{n_{\s (j)}}$, $M_{ij}=0 $ для иных $i$.


Это означает, что в столбце с номером $n_{\s (1)}+\dotsb +n_{\s (j-1)}+k$,


$1\leq k\leq n_{\s (j)}$,


проходящем через $j$-й блок, все элементы нулевые, кроме того элемента, который


находится в строке с номером $n_1+\dotsb +n_{\s (j)-1}+k$.





Покажем, что над числом $n_1+\dotsc +n_{\s (j)-1}+k$,


$1\leq k \leq n_{\s (j)}$, в подстановке


$\al * \s =


\begin{pmatrix}


{\bold p}_1 & {\bold p}_2 & \dotsc & {\bold p}_m\\


{\bold p}_{\s (1)} & {\bold p}_{\s (2)} & \dotsc & {\bold p}_{\s (m)}


\end{pmatrix}


$


расположено число $n_{\s (1)}+\dotsb +n_{\s (j-1)}+k$.


Число $n_1+\dotsb +n_{\s (j)-1}+k$ 


расположено в блоке ${\bold p}_{\s (j)}$ на $k$-м


месте. Число, расположенное над ним, есть его порядковый номер, считая слева


направо. Следовательно, это число элементов в блоках


${\bold p}_{\s (1)},\dotsc ,{\bold p}_{\s (j-1)}$


плюс $k$: $n_{\s (1)}+\dotsb +n_{\s (j-1)}+k$.


Рассмотрим


$$


M'=


\begin{pmatrix}


M'_{11} & \dotsc & M'_{1m} \\


\dotsc & \ddots & \dotsc \\


M'_{m1} & \dotsc & M'_{mm} 


\end{pmatrix}


=M^T=M^{-1},


$$


где $M'_{ij}$ --- блок размером $n_{\s (i)}\times n_j$, причем


$M'_{i,\s (i)}=E_{n_{\s (i)}}$, $M'_{ij}=0$ для иных $j$.


Вычислим $MCM'$.


Это матрица из $m\times m$ блоков, причем $(i,j)$-й блок имеет


размер $n_i\times n_j$, и вычисляется по формуле


$\sum\limits_{k=1}^m \sum\limits_{l=1}^m M_{ik}C_{kl}M'_{lj}$.


Ненулевое слагаемое в этом выражении получается лишь при


$k=\s ^{-1}(i)$, $l=\s ^{-1}(j)$.


В этом случае


$$


M_{i,\s ^{-1}(i)}=E_{n_{\s (\s ^{-1}(i))}}=E_{n_i},\quad


M'_{\s ^{-1}(j),j}=E_{n_{\s (\s ^{-1}(j))}}=E_{n_j},


$$


и вся сумма равна $C_{\s ^{-1}(i),\s ^{-1}(j)}$.





Итак, $\al * \s $ действует на $C$ следующим образом:


$$


(\al * \s )C=


\begin{pmatrix}


C_{\s ^{-1}(1),\s ^{-1}(1)} & \dotsc & C_{\s ^{-1}(1),\s ^{-1}(m)} \\


\dotsc & \ddots & \dotsc \\


C_{\s ^{-1}(m),\s ^{-1}(1)} & \dotsc & C_{\s ^{-1}(m),\s ^{-1}(m)}


\end{pmatrix}.


$$


На $(i,j)$-м месте в этой матрице стоит блок $C_{\s ^{-1}(i),\s ^{-1}(j)}$.


Если $i=j$, то


$C_{\s ^{-1}(i),\s ^{-1}(i)}=A_{\s (\s ^{-1}(i))}=A_i$;


eсли $i\ne j$, $k=\s ^{-1}(i)$, $l=\s ^{-1}(j)$, то


$C_{kl}=b_{kl}J_{n_{\s (k)},n_{\s (l)}}=b_{\s ^{-1}(i),\s ^{-1}(j)}J_{n_i,n_j}$.





Отсюда следует $(\s * \al)C=(A_1\dotsc A_m)(\s B)$.


\end{pf}








Пусть $\J (n)$ --- множество всех (не обязательно простых) 


неориентированных конечных графов с $n$ вершинами, помеченными 


символами $1, 2,\dotsc ,n$. Будем считать, что графы рассматриваются


с точностью до изоморфизма следующего вида: графы


$\G ', \G ''\in\J (n)$ считаются равными, 


если между множествами ребер, соединяющих любые две вершины с метками 


$i$ и $j$ (в дальнейшем будем отождествлять вершины и их метки) 


в графе $\G '$, и ребер, соединяющих вершины $i$ и $j$ 


в графе $\G ''$, можно установить 


взаимно однозначное соответствие. 





Определим на множестве $\J (n)$ структуру операды.


Как обычно, для графа $\G \in\J (n)$ через $ V(\G )$ обозначается


множество его вершин, а через $E(\G )$ --- множество его ребер.


Зададим композицию


$$


\J(n_1)\times \dotsb \times \J(n_m)\times\J(m)\rightarrow \J(n_1+\dotsb+n_m),


\quad (\G_1,\dotsc,\G_m,\G_0)\mapsto \G_1\dotsc\G_m\G_0.


$$


Пусть $\G_i\in\J(n_i)$, $1\leq i\leq m$, $\G_0\in\J(m)$.


Граф $\G=\G_1\dotsc \G_m\G_0$ устроен следующим образом:


$V(\G)=\{ 1, 2, \dotsc , n_1, n_1+1, \dotsc , n_1+n_2, n_1+n_2+1, \dotsc ,


n_1+\dotsc +n_m\}$.


Множество


$\{ n_1+\dotsb +n_{i-1}+1, \dotsc , n_1+\dotsb +n_{i-1}+n_i \}$ будем 


отождествлять с $\{ 1, \dotsc , n_i \} =V(\G_i)$,


сопоставляя вершине $1\leq j\leq n_i$ вершину $n_1+\dotsc +n_{i-1}+j$.


Таким образом, можно считать, что $V(\G)=\bigcup\limits_{i=1}^m V(\G_i)$,


причем $V(\G_i)\bigcap V(\G_j)=\emptyset$. 





Ребра $E(\G)$ описываются так. Пусть ребро $e\in E(\G_i)$ 


соединяет вершины $u, v$, где $1\leq u, v\leq n_i$.


Тогда в графе $\G$ определено ребро с тем же именем $e$, 


соединяющее вершины $n_1+\dotsb +n_{i-1}+u$ и 


$n_1+\dotsb +n_{i-1}+v$. С учетом сделанного выше 


отождествления вершин $\G_i$ с подмножеством


вершин $\G$ это означает, что $E(\G_i)$ вкладывается в $E(\G)$ так, что граф


$\G_i$ --- это подграф графа $\G$ при $1\leq i\leq m$. 





Кроме того, если $e\in E(\G_0)$ соединяет вершины 


$i$ и $j$ ($1\leq i, j\leq m$) графа $\G_0$, то для любых вершин


$u\in E(\G_i)$, $v\in E(\G_j)$ определено ребро


$e_{uv}=e_{vu}\in E(\G)$, соединяющее вершину $n_1+\dotsb +n_{i-1}+u$


графа $\G$ с вершиной


$n_1+\dotsb +n_{j-1}+v$.





Определим действие $\sym_n$ на $\J (n)$. Пусть $\s\in\sym_n$,


$\G\in\J (n)$. Положим $\s\G$ равным графу с вершинами $1,\dotsc , n$, 


причем $i$ и $j$ соединены ребром $e$ в $\G$ тогда и только тогда, когда 


вершины $\s (i)$ и $\s (j)$ соединены ребром с тем же именем $e$ в графе $\s\G$.





Проверим, что так определено действие, т.\,е.


$(\s\t )\G=\s (\t\G)$. Пусть $\G_1=\t\G$, $G_2=\s (\t\G)$, $\G_3=(\s\t )\G$.


Надо показать, что $\G_2=\G_3$, т.\,е. согласно сделанным предположениям 


должно существовать взаимно однозначное соответствие между ребрами, 


соединяющими каждую пару вершин $u,\ v$ (вершины в $\G_2$ 


и $\G_3$ общие). Ребро $e$ соединяет в графе $\G$ вершины $i$ и $j$ 


тогда и только тогда, когда ребро с тем же  


именем $e$ соединяет $\s\t (i)$ и $\s\t (j)$ в $\G_3$.





С другой стороны, $e$ соединяет $i$ и $j$ в $\G$ 


тогда и только тогда, когда ребром с именем $e$ 


соединены вершины $\t (i)$ и $\t (j)$ в $\G_1$, 


что эквивалентно тому, что вершины 


$\s\t (i)$ и $\s\t (j)$ coединены ребром $e$. 


Поскольку других ребер, кроме ребер указанного вида, в $\G_2$ и $\G_3$


нет, то наличие взаимно однозначного соответствия очевидно.





Через $E$ обозначим граф с одной вершиной и пустым множеством ребер


($E\in\J (1)$).





\begin{thm}


$1)$ Семейство множеств $\J=\{\J (n) \mid n\geq 1\}$ c определенной 


на них выше операцией композиции и действием 


симметрических групп является операдой. Единица операды --- граф $E$.





$2)$


Подоперада операды $\J$, состоящая из графов без петель, изоморфна


подопераде, описанной в теореме $1$ операды $\M$ $($где в качестве $K$ берется


множество целых неотрицательных чисел и $\varepsilon=0)$, состоящей


из симметрических матриц.


\end{thm}





\begin{pf}


Проверим свойства единицы. Пусть $\G_1=\dotsb =\G_m=E$,


$\G=E\dotsc E\G_0$. В графе $\G$ имеется $m$ вершин, как


и в $\G_0$. Поскольку $\G_i=E$, $1\leq i\leq m$, 


не имеет ребер, то все ребра $\G$ получаются 


из ребер $\G_0$. Пусть $k, \ l$ --- вершины $\G_0$, $e$~--- соединяющее 


их ребро в $\G_0$. Вершины $\G_k$ и $\G_l$ можно выбрать единственным 


способом, и они имеют метки (в $\G$) $k$ и $l$. Таким образом, по определению 


композиции, в $\G$ найдется лишь одно ребро $e_{kl}$. Теперь очевидно, 


что имеется взаимно однозначное


соответствие между ребрами $\G$ и $\G_0$.


Следовательно, это один и тот же элемент $\J (m)$.


Если $m=1$, $\G_0=E$, $\G=\G_1\G_0=\G_1E$, то прямо из определения 


композиции ясно, что вершины и ребра $\G$ 


те же, что и у $\G_1$.





Докажем, что композиция ассоциативна.  


Пусть $\ov\G_i=\G_{1i}\dotsc \G_{n_ii}$,


$1\leq i\leq m$, $\G_{ji}\in \J (k_{ij})$, $\G_i\in \J (n_i)$, $\G_0\in \J (m)$.


Покажем, что графы 


$\G' =(\ov\G_1\G_1)\dotsc (\ov\G_m\G_m)\G_0$ и


$\G'' =\ov\G_1\dotsc \ov\G_m (\G_1\dotsc \G_m\G_0)$ 


представляют один и тот же элемент из 


$\J \Big(\sum\limits_{i=1}^m\sum\limits_{j=1}^{n_i}k_{ji}\Big)$.





Ясно, что множество вершин у этих графов одно и то же: 


$\Big\{ 1, 2, \dotsc , \sum\limits_{i=1}^m


\sum\limits_{j=1}^{n_i}k_{ji}\Big\}$. При этом вложения графов 


$\G_{ji}\rightarrow \G_{1i}\dotsc \G_{n_ii}\G_i\rightarrow \G''$,


$\G_{ji}\rightarrow \G''$


отображают вершины $\G_{ji}$ в $V(\G')=V(\G'')$ одинаковым образом:


$l\mapsto \sum\limits_{s=1}^{i-1}\sum\limits_{t=1}^{n_s}k_{ts}+


\sum\limits_{t=1}^{n_i-1}k_{ti}+l$.





Отождествляя $\G_{ji}$ с подграфами $\G'$ и $\G''$, 


а $\G_i$ --- с подграфами $\G_1\dotsc \G_m\G_0$,


рассмотрим ребра $\G'$ и $\G''$. Множество ребер $\G'$ 


можно разбить на три непересекающихся подмножества:


\begin{itemize}


\item[1)]


$\bigcup \limits_{i=1}^m \bigcup \limits_{j=1}^{n_i} E(\G_{ji});$


\item[2)]


множество ребер, определяемых так: для каждого ребра $e\in E(\G_i)$, 


инцидентного вершинам $k, l\in V(\G_i)$,


для любых вершин $u\in V(\G_{ki})$, $v\in V(\G_{li})$ 


определено ребро $e_{uv}=e_{vu}$, cоединяющее 


$u$ и $v, 1\leq i\leq m;$ 


\item[3)]


множество ребер, определяемых так: для каждого ребра $e\in E(\G_0)$, 


соединяющего вершины $k,\ l$ графа 


$\G_0$, произвольных $p, \ q$, $1\leq p\leq n_k$,


$1\leq q\leq n_l$, и для любых $u\in V(\G_{pk})$, $v\in V(\G_{ql})$ 


определено ребро $e_{pq, uv}$, cоединяющее $u$ и $v$ 


(здесь $e_{pq, uv}=e_{qp, uv}=e_{pq, vu}$ и т.\,д.).


\end{itemize}





Аналогичным образом, используя определение композиции, 


множество ребер $\G''$ можно разбить на три 


непересекающихся подмножества, которые (с точностью до обозначений) 


совпадают с множествами из пп.\,1), 2), 3) выше. 


Поэтому $\G'=\G''$ как элементы операды $\J$.





Проверим свойства операды, связанные с подстановками. Пусть 


$\s_i\in \sym_{n_i}$, $\G_i\in \J (n_i)$,\linebreak $1\leq i\leq m$,


$\G_0\in \J (m)$; $\G'=(\s_1\G_1)\dotsc (\s_m\G_m)\G_0$, 


$\G''=(\s_1\times \dotsb \times \s_m)(\G_1\dotsc \G_m\G_0)$. 


Покажем, что это один и тот же элемент $\J (n_1+\dotsb +n_m)$. 


Для этого построим взаимно однозначное соответствие ребер 


(совпадение вершин очевидно). Пусть $e$ --- ребро $\G'$, соединяющее две 


вершины подграфа $\s_i\G_i\subset \G'$, например,


$n_1+\dotsb +n_{i-1}+\s_i(k)$ 


и $n_1+\dotsb +n_{i-1}+\s_i(l)$, $1\leq k, l\leq n_i$. 


Оно соответствует ребру $e$, соединяющему вершины 


$n_1+\dotsb +n_{i-1}+k$ и $n_1+\dotsb +n_{i-1}+l$


подграфа $\G_i$ графа $\G_1\dotsc \G_m\G_0$.   


Но по определению $\s_1\times \dotsb \times \s_m$ 


этому ребру соответствует ребро, соединяющее вершины


$(\s_1\times \dotsb \times \s_m)(n_1+\dotsb +n_{i-1}+k)=


n_1+\dotsb +n_{i-1}+\s_i(k)$ и 


$(\s_1\times \dotsb \times \s_m)(n_1+\dotsb +n_{i-1}+l)=


n_1+\dotsb +n_{i-1}+\s_i(l)$


графа $\G''=(\s_1\times \dotsb \times \s_m)(\G_1\dotsc \G_m\G_0)$.


Если же $e$ --- ребро $\G_0$, соединяющее вершины $k$ и $l$,


$1\leq k, l\leq m$, то любые вершины в $\s_k\G_k$ 


и $\s_l\G_l$ можно представить в виде $\s_ku$ 


и $\s_lv$, $1\leq u\leq n_k$, $1\leq v\leq n_l$, 


так что в $\G'$ определено ребро $e_{\s_ku, \s_lv}$, 


соединяющее вершины $n_1+\dotsb +n_{k-1}+\s_ku$ и


$n_1+\dotsb +n_{l-1}+\s_lv$. 


Для того же ребра $e$ в графе $\G_1\dotsc \G_m\G_0$ 


определено ребро $e_{u,v}$, соединяющее вершины 


$n_1+\dotsb +n_{k-1}+u$ и $n_1+\dotsb +n_{l-1}+v$.


Этому ребру соответствует в графе 


$(\s_1\times \dotsb \times \s_m)(\G_1\dotsc \G_m\G_0)$ ребро, 


соединяющее вершины 


$(\s_1\times \dotsb \times \s_m)(n_1+\dotsb +n_{k-1}+u)=


n_1+\dotsb +n_{k-1}+\s_ku$ и


$(\s_1\times \dotsb \times \s_m)(n_1+\dotsb +n_{l-1}+v)=


n_1+\dotsb +n_{l-1}+\s_lv$.





Пусть теперь $\s\in\sym_m$, $\G'=(\G_1\dotsc \G_m)(\s\G_0)$,


$\al=(n_1,\dotsc ,n_m)$, 


$\G''=(\al * \s)(\G_{\s 1}\dotsc \G_{\s m}\G_0)$.


Построим взаимно однозначное соответствие между ребрами $\G'$ и $\G''$. 





Пусть $e$ --- ребро графа $\G_{\s j}$, соединяющее вершины $k$ и $l$,


$1\leq k,l\leq n_{\s (j)}$, $1\leq j\leq m$. 


При вложении $\G_{\s j}$ в $\G'$ ему соответствует ребро 


(с тем же именем), соединяющее вершины 


$n_1+\dotsb +n_{\s (j)-1}+k$ и $n_1+\dotsb +n_{\s (j)-1}+l$. 





Рассмотрим вложение графа $\G_{\s (j)}$ в граф $\G_{\s 1}\dotsc \G_{\s m}\G_0$. 


Ребру $e$ графа $\G_{\s (j)}$ будет соответствовать ребро 


с тем же именем, соединяющее вершины 


$n_{\s (1)}+\dotsb +n_{\s (j-1)}+k$ и $n_{\s (1)}+\dotsb +n_{\s (j-1)}+l$.


Тогда в графе $(\al * \s)(\G_{\s 1}\dotsc \G_{\s m}\G_0)$ ребро с именем $e$


должно соединять вершины


$(\al * \s)(n_{\s (1)}+\dotsb +n_{\s (j-1)}+k)$ и


$(\al * \s)(n_{\s (1)}+\dotsb +n_{\s (j-1)}+l)$.


Но при доказательстве теоремы 1 уже было показано, 


что эти числа равны  


$n_1+\dotsb +n_{\s (j)-1}+k$ и $n_1+\dotsb +n_{\s (j)-1}+l$ соответственно.





Пусть $e$ --- ребро графа $\G_0$, соединяющее вершины 


$k$ и $l$, $1\leq k,l\leq m$.


Ему соответствует ребро $e$ графа $\s\G$, соединяющее $\s k$ и $\s l$.


Тогда для любых вершин $u\in V(\G_{\s (k)})$,


$1{\leq} u{\leq} n_{\s (k)}$, $v\in V(\G_{\s (l)})$, $1\leq v\leq n_{\s (l)}$, 


определено ребро $e_{u,v}$ в $\G'$, соединяющее 


$n_1+\dotsb +n_{\s (k)-1}+u$ и $n_1+\dotsb +n_{\s (l)-1}+v$.  


С другой стороны, по ребру $e\in E(\G_0)$ и вершинам 


$u\in V(\G_{\s (k)})$,\linebreak $v\in V(\G_{\s (l)})$


строится ребро $e_{u,v}$ в графе $\G_{\s 1}\dotsc \G_{\s m}\G_0$, 


соединяющее вершины  


$n_{\s (1)}+\dotsb +n_{\s (k-1)}+u$ и $n_{\s (1)}+\dotsb +n_{\s (l-1)}+v$.  


В графе $(\al * \s)(\G_{\s 1}\dotsc \G_{\s m}\G_0)$ этому 


ребру соответствует ребро с тем же именем, соединяющее вершины 


$(\al * \s)(n_{\s (1)}+\dotsb +n_{\s (k-1)}+u)=n_1+\dotsb +n_{\s (k)-1}+u$ и


$(\al * \s)(n_{\s (1)}+\dotsb +n_{\s (l-1)}+v)=n_1+\dotsb +n_{\s (l)-1}+v$.





Тем самым проверены все свойства операды для $\J$. 





Пусть $K=\{0, 1, 2,\dotsc \}$ --- множество 


целых неотрицательных чисел. Каждому графу $\G\in\J (n)$ можно


сопоставить матрицу инциденций $A=A(\G )$, компоненты которой ---


элементы~$K$.


Так как вершины графов из $\J (n)$ упорядочены, то таким образом задается


инъективное отображение из


$\J (n)$ в $\M_{n,n}=\ov\M (n)$.


Множество графов из


$\J (n)$ без петель при этом отображается в $\M (n)$.





Легко убедиться, что если $\G_1, \dotsc , \G_m, \G_0$ --- графы 


без петель, то и $\G_1\dotsc \G_m\G_0$ обладает этим свойством, 


и если матрицы $A_1\in\M (n_1),\dotsc , A_m\in\M (n_m)$, $A_0\in\M (m)$


симметричны, то и матрица $A_1\dotsc A_mA_0\in\M (n_1+\dotsc +n_m)$ 


также симметрична. Таким образом, эти графы и матрицы образуют 


подоперады в $\J$ и $\M$ соответственно.





Матрица инциденций для неориентированного графа является 


симметричной, ее элементы --- целые неотрицательные числа, 


и т.\,к. рассматриваются графы без петель, то диагональные 


элементы матрицы равны нулю. Обратно, любая матрица такого 


вида однозначно определяет граф без петель --- элемент $\J (n)$. 


Легко проверяется, что $A(\G_1\dotsc \G_m\G_0)=A(\G_1)\dotsc A(\G_m)A(\G_0)$, и


$A(\s \G)=\s A(\G)$ при $\s\in \sym_n$, $\G\in \J (n)$.


\end{pf}








\begin{note}


Аналогично тому, как это сделано выше для неориентированных графов, 


можно определить операду, элементами $n$-й компоненты которой будут 


ориентированные графы с множеством вершин $\{ 1, 2,\dotsc ,n\}$.


Можно также ввести операды взвешенных графов (ориентированных или нет), 


ребрам которых приписаны веса из моноида $K$. Имеется также несколько 


структур операд на множествах помеченных гиперграфов. В данной работе, однако, 


сосредоточимся на операде неориентированных графов. 


\end{note}











\begin{defn}


Граф $\G$ будем называть {\it операдно разложимым} (или просто разложимым), 


если существует такая нумерация его


вершин, что $\G$ изоморфен операдной композиции 


$\G_1\dotsc \G_m\G_0$, где $\G_1,\dotsc ,\G_m,\G_0$ --- помеченные графы


(элементы операды $\J$) такие, что $\G_0\ne E$ 


и, по крайней мере, один из $\G_i$


нетривиален при $i\geq1$.





Граф $\G$ будем называть {\it операдно неразложимым} (или просто неразложимым), 


если он не является операдно разложимым.


Иными словами, для неразложимого графа


из $\G\cong\G_1\dotsc \G_m\G_0$ следует, что либо


$\G_1=\dotsb =\G_m=E$, $\G_0\cong\G$,


либо $m=1$, $\G_0=E$, $\G_1\cong\G$.


\end{defn}





\begin{lem}


Пусть $\G=\G_1\dotsc \G_m\G_0$. 


Тогда в $\G$ существует подграф, изоморфный $\G_0$.


\end{lem}





\begin{pf}


Строим этот подграф следующим обрaзом. Выбираем для каждого\linebreak


$i\in V(\G_0)$


по одной вершине $v_i\in V(\G_i)$, $1\leq i\leq m$, и для каждого ребра $e$, 


соединяющего в $\G_0$ вершины $i$ и $j$ --- ребро $e_{v_iv_j}$, соединяющее


$v_i\in V(\G_i)\subset V(\G)$ и $v_j\in V(\G_j)\subset V(\G)$. 


Очевидно, подграф $\G'$ графа $\G$ c множеством вершин 


$\{ v_1,\dotsc ,v_m\}$ и множеством ребер $\{ e_{v_iv_j} \mid e\in E(\G_0)\}$ 


изоморфен $\G_0$.


\end{pf}





\begin{defn}


{\it Фрагментом} графа $\G$ с ядром $\G^*$


будем называть пару графов $(\G', \G^*)$, где $\G^*\subseteq \G'\subseteq \G$,


причем выполнены условия


\begin{itemize}


\item[1)]


eсли $v_1, v_2\in V(\G^*)$ инцидентны ребру $e\in E(\G)$, то $e\in E(\G^*)$;


\item[2)]


eсли $v\in V(\G^*)$ соединено ребром $e$ c $v'\in V(\G)$,


то $v'\in V(\G')$ и $e\in E(\G')$.


\end{itemize}





Фрагмент $(\G', \G^*)$ будем называть {\it нетривиальным}, если ядро $\G^*$


не равно тривиальному графу с одной вершиной без ребер, и


$V(\G')\ne V(\G^*)$.





Нетривиальный фрагмент $(\G', \G^*)$ графа $\G$ 


будем называть {\it разлагающим}, если выполняется еще одно условие


\begin{itemize}


\item[3)]


eсли вершины $v\in V(\G^*)$ и $v'\in V(\G')$, $v'\notin V(\G^*)$, 


соединяет (в графе $\G$)


$k$ ребер ($k\geq 0$ согласно условию~2) 


все эти ребра принадлежат графу $\G'$), то любая другая


вершина $w\in V(\G^*)$ соединена в графе $\G$ c $v'$ 


в точности $k$ ребрами.


\end{itemize}


\end{defn}





\begin{note}


Если $(\G', \G^*)$ --- фрагмент, то для 


$\G\supseteq \G''\supseteq \G'$ пара $(\G'', \G^*)$ также будет 


фрагментом $\G$, разлагающим, если разлагающим был фрагмент 


$(\G', \G^*)$. 


\end{note}





\begin{lem}


Фрагмент $(\G', \G^*)$ является разлагающим тогда и только тогда, 


когда после некоторой перенумерации вершин $\G'\cong\G^*E\dotsc E\G_0$ 


для некоторого графа $\G_0$, в котором 


нет петель, инцидентных вершине $1$. 


\end{lem}





\begin{pf}


Пусть $\G'=\G^*E\dotsc E\G_0$ и $1, 2, \dotsc , m$ --- вершины $\G_0$.


Положим $\G_1=\G^*$, $\G_2=E, \dotsc ,\G_m=E$,


так что $\G'=\G_1\G_2\dotsc\G_m\G_0$. Тогда $\G_1=\G^*$ --- подграф $\G'$.


Пусть $1,\dotsc, n_1$ --- вершины $\G_1$, тогда остальными вершинами $\G'$ будут  


$n_1+1, \dotsc , n_1+m-1$. Эти вершины соответствуют тривиальным 


подграфам $E$ в разложении


$\G'=\G^*E\dotsc E\G_0$. По определению операдной композиции каждому ребру $e$


графа $\G_0$, соединяющему вершины $1$ и $j>1$, соответствует 


семейство ребер $\G'$ вида $e_{i, n_1+j}$, соединяющих все вершины $\G_1=\G^*$


(т.\,е. вершины $ i=1, \dotsc ,n_1$) с вершиной $n_1+j$. 


Если в $\G_0$ имеется $k$ ребер, инцидентных $1$ и $j$, то для $\G'$ 


это будет означать, что каждая вершина $\G^*$ соединена с вершиной $n_1+j$ 


одним и тем же количеством $k$ ребер.


Таким образом, условие~3) из определения выполнено.





Обратно, пусть $(\G', \G^*)$ --- разлагающий фрагмент,


$n_1=\vert{ V(\G^*)\vert}$.


Снабдим вершины $\G'$ метками, начав нумерацию с вершин $\G^*$, так что 


$V(\G')=$ $\{ 1, 2, \dotsc ,n_1$, $n_1+1, \dotsc , n_1+m-1\} $ 


и $V(\G^*)=\{ 1, \dotsc , n_1\}$. Пусть $\G_0$ --- граф, полученный 


``стягиванием'' всех ребер и вершин $\G^*$ в одну точку. Множеством его 


вершин будет $\{ 1, 2, \dotsc , m\}$, причем вершина~$1$ cоответствует 


всем вершинам $\G^*$, вершины $j$ --- вершинам $n_1+j$ при $2\leq j\leq m-1$.


Множество ребер $\G_0$ устроено следующим образом. Если $e$ --- ребро $\G'$, 


соединяющее $j$ и $t$, где $j, t>n_1$, то ему соответствует 


ребро $\G_0$, соединяющее $j-n_1$ и $t-n_1$.  


Для любого $j>n_1$ либо не существует ребер $\G'$, cоединяющих $j$ с вершинами 


$1, \dotsc , n_1$, либо каждая вершина $1, \dotsc , n_1$ соединена с $j$ 


одним и тем же количеством (напр., $k$) ребер. В этом случае в $\G_0$ 


вершина $1$ соединена с вершиной $j-n_1$ ровно $k$ ребрами.


Пусть $\G_1=\G^*$, $\G_2=\dotsb =\G_m=E$. Тогда, сравнивая 


$\G'$ c $\G_1\dotsc \G_m\G_0$,


прямо из определения делаем вывод об изоморфности этих графов.


\end{pf}








\begin{lem}


Пусть $\G'=\G_1E\dotsc E\G'_0$ для некоторого графа 


$\G'_0$. Тогда существуют графы 


$\G_0$ и $\G^*$ такие, что $\G'=\G^*E\dotsc E\G_0$, 


и в $\G_0$ нет петель, инцидентных вершине $1$. 


\end{lem}





\begin{pf}


Если в $\G'_0$ нет петель, инцидентных вершине $1$, 


то $\G_0=\G'_0$, $\G^*=\G_1$, и все доказано. В противном случае, 


пусть $\Lambda $ --- подграф $\G'_0$ с единственной 


вершиной~$1$, ребра которого --- всевозможные петли, инцидентные этой вершине.





Пусть $\G_0$ --- граф, получаемый из $\G'_0$ стягиванием всех ребер $\Lambda$.


Множество вершин у $\G_0$ то же, что и у $\G'_0$, но в $\G_0$ уже нет петель,


инцидентных вершине~$1$. 


Кроме того, очевидно, что $\G'_0=\Lambda E\dotsc E\G'_0$


(где тривиальные графы $E$ соответствуют вершинам $\G'_0$, 


не входящим в $\Lambda $). Отсюда


$\G_1E\dotsc E\G'_0=\G_1E\dotsc E(\Lambda E\dotsc E\G_0)=


(\G_1\Lambda )(EE)\dotsc (EE)\G_0=


(\G_1\Lambda )E\dotsc E\G_0$, и можно взять $\G^*=\G_1\Lambda $.


\end{pf}





\begin{lem}


Пусть $\G=\G_1\dotsc \G_{n_1}\G_{n_1+1}\dotsc \G_m\G_0$ 


и выполнены следующие условия\/${:}$


\begin{itemize}


\item[a)]


$\G_1, \dotsc ,\G_{n_1}$ --- пустые графы, т.\,е. графы без ребер$;$


\item[b)] в $\G_0$ имеется фрагмент $(\G'_0, \G^*_0)$ такой, 


что $V(\G^*_0)=\{ 1, \dotsc , n_1\}$.


\end{itemize}


Тогда в $\G$ можно выбрать фрагмент $(\G', \G^*)$ такой, что $\G^*\cong\G^*_0$.





Если к тому же $(\G'_0, \G^*_0)$ --- разлагающий фрагмент, то и $(\G', \G^*)$


можно выбрать разлагающим.


\end{lem}  





\begin{pf}


Воспользуемся леммой 3, в которой строится подграф $\G''\subset \G$, 


изоморфный $\G_0$. Пусть $\G^*$ --- подграф, являющийся образом $\G^*_0$ 


при этом изоморфизме, а $\G'$ строится так, что 


$V(\G')=\bigcup\limits_{j\in V(\G'_0)}V(\G_j)$, и если $v_1\in V(\G_{j_1})$, 


$v_2\in V(\G_{j_2})$, то ребра $V(\G')$, cоединяющие $v_1$ и $v_2$, имеют 


вид $e_{v_1v_2}$, где $e\in V(\G')$ соединяет $j_1$ и $j_2$.


Проверим условия~1) и~2) из определения фрагмента. Пусть $v_1, v_2\in V(\G^*)$.


Это значит, что $v_1\in V(\G_{i_1})$, $v_2\in V(\G_{i_2})$,


$i_1, \ i_2$ --- вершины $\G_0$ и


вершины $\G^*_0$ соответственно, так что $1\leq i_1, i_2\leq n_1$. 


Пусть $e$ --- ребро $\G$, инцидентное $v_1$ и $v_2$. Допустим, что 


$v_1=v_2=v$. Тогда по построению $\G''$ должно быть 


$i_1=i_2=i$. Так как в $\G_i$ нет петель, то $e=u_{vv}$, 


где $u$ --- петля в $\G_0$, инцидентная вершине $i$.


Так как $i\in V(\G^*_0)$, то из условия~1) 


для фрагмента $(\G'_0, \G^*_0)$ следует $u\in E(\G^*_0)$. Отсюда 


по определению $\G''$ будем иметь $u_{vv}\in E(\G^*)$.


Если же $v_1\ne v_2$, то в этом случае $i_1\ne i_2$, 


и ребро $e$ автоматически имеет вид


$u_{v_1v_2}$ для некоторого ребра $u\in V(\G_0)$, инцидентного $i_1$ и $i_2$.


Но это ребро обязано принадлежать $E(\G^*_0)$, 


и по построению $u_{v_1v_2}\in E(\G^*)$.





Пусть теперь даны вершины $v\in V(\G^*)$, $v'\in V(\G)$ 


и соединяющее их ребро $e\in E(\G)$.


Условие~a) исключает случай, когда $v$ и $v'$ 


принадлежат одному и тому же $\G_i$, $1\leq i\leq m$.


Пусть $v\in V(\G_i)$, $v'\in V(\G_j)$, $i\ne j$. 


Тогда $e=u_{vv'}$, где $u$ --- ребро $\G_0$, cоединяющее


вершины $i\in V(\G^*_0)$ и~$j$. По определению фрагмента 


должно быть $j\in V(\G'_0)$ и $u\in E(\G'_0)$.


Это означает, что $v'\in V(\G')$ и $e=u_{vv'}\in E(\G')$. 





Пусть фрагмент $(\G'_0, \G^*_0)$ разлагающий. 


Проверим условие~3) для фрагмента $(\G', \G^*)$.


Пусть $v\in V(\G^*)$, $v\in V(\G_i)$, 


$1\leq i\leq n_1$, $v'\in V(\G')$, $v'\in V(\G_j)$, $j>n_1$, 


и имеется ровно $k$ ребер


$e^{(1)}, \dotsc , e^{(k)}$, соединяющих в $\G$ вершины $v$ и $v'$.


Так как $j\ne i$, то все эти ребра имеют вид $e^{(t)}=u_{ij}^{(t)}$,


где $u^{(1)}, \dotsc , u^{(k)}$ --- ребра $\G_0$,


cоединяющие вершины $i$ и $j$.


Пусть дана другая вершина $u\in V(\G^*)$, $u\in V(\G_p)$, $1\leq p\leq n_1$.


В $\G_0$ имеется ровно $k$ ребер, соединяющих вершину $j$ с вершиной $p$.


Если $w^{(1)}, \dotsc , w^{(k)}$ --- эти ребра, то им соответствуют в


$\G$ ровно $k$ ребер


$w^{(1)}_{i, p}, \dotsc , w^{(k)}_{i, p}$, соединяющих вершины $u$ и $v'$.


По определению операдной композиции графов других ребер, соединяющих $u$


и $v'$, быть не может.


\end{pf}





\begin{thm}


Граф $\G\in \J (n)$ разложим тогда и только тогда, когда он


содержит нетривиальный разлагающий фрагмент.


\end{thm}





\begin{pf}


Пусть $\G=\G_1\dotsc\G_m\G_0$ --- нетривиальное разложение. Можно считать, что


$\G_1\ne E$, и пусть $n_1=| V(\G_1) |$. Тогда


$$


\G=\G_1\dotsc\G_m\G_0=(\underbrace{E\dotsc E}_{n_1}


\G_2\dotsc\G_m)(\G_1\underbrace{E\dotsc E}_{m-1}\G_0). 


$$





По доказанному выше можно считать, что 


$\G'_0=\G_1E\dotsc E\G_0=\G^*_0E\dotsc E\G''_0$,


где $\G''_0$ --- граф без петель, инцидентных вершине $1$, 


и тогда $(\G'_0, \G^*_0)$ --- разлагающий фрагмент в $\G'_0$. 


Из леммы~6 делаем вывод о том, что разлагающий фрагмент 


существует и в $\G$. 





Обратно, пусть в $\G$ существует нетривиальный разлагающий фрагмент


$(\G', \G^*)$.


Исходя из замечания к определению 2, можно считать, 


что $\G'=\G$. Тогда из леммы 4 следует


$\G\cong\G^*E\dotsc E\G_0$, т.\,е. $\G$ разложим.


\end{pf} 





\begin{note}


Несвязный граф $\Gamma$ с более чем двумя вершинами


разложим, т.\,к. любое объединение связных компонент, содержащее $k$ вершин,


$1< k < |V(\Gamma)|$, будет ядром разлагающего фрагмента. В частности,


операдно неразложимые графы с более чем двумя вершинами связны.


\end{note}








Аналогично тому, как это сделано выше, определяется операда простых 


помеченных графов без петель $G=\{ G(n) \mid n\geq 1\} $. Определение 


разлагающего фрагмента, его свойства, формулировка и доказательство 


теоремы 3 для случая простых графов остаются теми же 


самыми. Ниже рассматриваются только простые графы без петель.


Заметим еще, что основные результаты работы \cite{LM} на языке операд


означают, что некоторые классы простых графов являются подоперадами


в операде $G$.





\begin{thm} 


Простой граф $\G$ с не менее чем тремя вершинами операдно 


неразложим тогда и 


только тогда, когда он удовлетворяет следующему свойству\/${:}$


для любого $2\leq n< | V(\G) |$ и любых $n$ вершин $v_1,\dotsc , v_n$ 


найдется вершина $v\notin \{ v_1,\dotsc , v_n\}$, 


и вершины $v_i$ и $v_j$, $i\neq j$, такие, что $v$ соединена ребром с


$v_i$, но не соединена ребром с $v_j$.


\end{thm}





\begin{pf}


Утверждение теоремы логически равносильно утверждению теоремы 3 для 


простых графов. Точнее, утверждение ``$\G$ разложим $\Leftrightarrow$ 


существует разлагающий фрагмент $(\G, \G^*)$'' логически равносильно 


утверждению ``$\G$ неразложим $\Leftrightarrow$ для любого $\G^*$ 


фрагмент $(\G, \G^*)$ не является разлагающим'', 


что сводится к тому , что для любого $\G^*$ нарушено свойство 3 из 


определения разлагающего фрагмента. В формулировке этого свойства 


участвуют только вершины $v_1,\dotsc , v_n$ подграфа $\G^*$, и его 


невыполнение означает, что найдется вершина $v\notin V(\G^*)$, 


которая соединена ребрами не со всеми вершинами $\G^*$


(с некоторыми соединена, с некоторыми нет).


\end{pf}





\begin{thm}


Пусть $\G$ --- простой связный граф без петель. Допустим, что $\G$ не 


содержит подграфов, изоморфных $K_3$ $($треугольников$)$, а также 


подграфов, изоморфных $K_{n, m}$ $(n\geq 2)$, устроенных следующим 


образом. Это подграфы с вершинами $v_1,\dotsc , v_n$, $u_1,\dotsc , u_m$,


причем каждая вершина $v_i$ соединена ребром с каждой вершиной $u_j$, 


а вершины $v_i$ могут быть соединены в $\G$ ребрами, кроме вершин вида 


$u_j$, только с вершинами вида $v_k$. 


Тогда граф $\G$ операдно неразложим.


\end{thm} 





\begin{pf}


Рассмотрим разложимый граф $\G$ (простой, связный, без петель), и пусть 


$\G^*$ --- ядро разлагающего фрагмента, $v_1,\dotsc , v_n$ ($n\geq 2$) 


--- вершины $\G^*$. Пусть $u_1,\dotsc , u_m$ --- все вершины $\G$, 


соединенные ребрами с вершинами $\G^*$. Ввиду связности $\G$ имеем 


$m>0$. Если $\G^*$ не дискретен, то две его вершины $v_i$ и $v_k$,


соединенные ребром, вместе с вершиной $u_j$ образуют подграф, изоморфный $K_3$.


Если даже $\G^*$ дискретен, то подграф $\G$ с множеством вершин


$v_1,\dotsc , v_n$, $u_1,\dotsc , u_m$ и множеством ребер, соединяющих 


вершины $v_i$ с вершинами $u_j$, изоморфен $K_{n, m}$, и устроен так, 


как это описано в формулировке теоремы. Если же в $\G$ нет таких подграфов, 


то не существует и нетривиальных разлагающих фрагментов.


\end{pf}





Применим теорему 5 для решения вопроса о неразложимости двух семейств графов.





\begin{exam}


Рассмотрим семейство графов --- многомерных кубов $Q_n$. Вершины $Q_n$ --- 


последовательности нулей и единиц длины $n$,


$\ov{x}=(x_1,\dotsc , x_n)$, $x_i=0$ или $x_i=1$,


$\vert{V(Q_n)}\vert=2^n$.


Две вершины $\ov x$ и $\ov y$ считаются соединенными ребром, если расстояние 


Хэмминга $d(\ov x,\ov y)=|\{ i \mid x_i\neq y_i\} |$ между ними равно единице. 


Покажем, что при $n>2$ граф $Q_n$ операдно неразложим. Граф $Q_2$ разложим 


следующим образом: $Q_2=\G_1\G_2K_2$, где $\G_1=\G_2$ --- дискретные графы 


с двумя вершинами. Покажем, что при $n>2$ в $Q_n$ нет ни треугольников, ни 


подграфов $K_{n, m}$ такого вида, который описан в теореме 5. 


Пусть в $Q_n$ существует треугольник с вершинами


$\ov{x}=(x_1,\dotsc , x_n)$, $\ov{y}=(y_1,\dotsc , y_n)$,


$\ov{z}=(z_1,\dotsc , z_n)$,


$d(\ov{x},\ov{y})=d(\ov{y},\ov{z})=d(\ov{x},\ov{z})=1$.


Допустим, что $x_1\ne y_1$, $x_2=y_2,\dotsc , x_n=y_n$.


Для $\ov{z}$ рассмотрим два случая.


Если $z_1\ne x_1$, то из $d(\ov{x},\ov{z})=1$ следует 


$z_2=x_2=y_2,\dotsc , z_n=x_n=y_n$.


Но из $x_1\ne y_1$, $x_1\ne z_1$ следует $y_1=z_1$ и $\ov{y}=\ov{z}$.


Если же $z_1=x_1$, и, например, $z_2\ne x_2=y_2$, 


$z_3=x_3=y_3,\dotsc , z_n=x_n=y_n$, то получим


$z_1=x_1\neq y_1$, и $z_2\ne y_2$, $d(\ov{z},\ov{y})=2$, что невозможно.


Степень каждой вершины $Q_n$ равна $n$. 


Пусть $\ov x$, $\ov z$ --- две различные вершины,


имеющие один и тот же набор смежных вершин $\ov {y_1},\dotsc , \ov {y_n}$.


Покажем, что этого не может быть.


Можно считать, что $\ov{y_i}=(x_1,\dotsc ,\wt {x_i},\dotsc , x_n)$,


где $\wt {x_i}=0$, если $x_i=1$, и $\wt {x_i}=1$, если $x_i=0$.


Тогда существует неединичная подстановка $\t\in\sym_n$ такая, что


$\ov y_{\t (i)}=(z_1,\dotsc ,\wt {z_i},\dotsc , z_n)$. 


Выберем $i$ так, чтобы $\t (i)\neq i$. 


Тогда $\ov x$ и $\ov y_{\t (i)}$ отличаются в $\t (i)$-й


компоненте (и только в ней), $\ov y_{\t (i)}$ и $\ov z$ 


отличаются только в $i$-й компоненте.


Так как $\t (i)\neq i$, $\ov x$ и $\ov z$ отличаются друг 


от друга и в $i$-й и в $\t (i)$-й компонентах.


Допустим, что существует $j\neq i$, $j\neq \t (i)$ такое, что $j\neq \t (j)$.


Тогда, рассуждая аналогичным образом, приходим к 


выводу, что $\ov x$ и $\ov z$ различаются


в $j$-й и $\t (j)$-й компонентах. 


Не исключается, что $\t (j)=j$, но, по крайней мере, 


$d(\ov x, \ov z)\geq 3$. Однако из


$d(\ov{x},\ov{y_i})+d(\ov{y_i},\ov{z})\geq d(\ov{x},\ov{z})\geq 3$


следует, что одно из чисел $d(\ov{x},\ov{y_i}), d(\ov{y_i},\ov{z})$


строго больше единицы, --- противоречие.


Допустим теперь, что $\t$ --- транспозиция, 


$\t (i)=k$, $\t (k)=i$, и $\t (j)=j$ при $j\neq i, k$.


При $n\geq 3$ такое число $j$ существует.


Сравнивая $\ov x, \ov y_i=\ov y_{\t (k)}$,


$\ov y_k=\ov y_{\t (i)}$, $\ov y_j=\ov y_{\t (j)}$ с $\ov z$,


заключаем, что $\ov z$ и $\ov y_j$ 


различаются сразу в трех компонентах: $i$-й, $j$-й и $k$-й. Полученное 


противоречие показывает, что в $Q_n$ при $n>3$ не существует подграфа, 


изоморфного $K_{p, q}$, удовлетворяющего условию из теоремы 5. 


\end{exam}





\begin{exam}


Опишем разлагающие фрагменты в деревьях. Будем называть вершину дерева~$T$ 


внешней, если ее степень равна единице, и внутренней в противном случае.


Гроздью в дереве~$T$ назовем поддерево $T'$ c вершинами 


$v_1,\dotsc , v_k, v_0$, $k\geq2$,


где $v_1,\dotsc , v_k$ --- внешние для $T$, а $v_0$ --- внутренняя для


$T$ вершина.





Утверждается, что наличие гроздьев в дереве равносильно наличию разлагающих 


фрагментов. Точнее, если $T'$ --- гроздь, $T^*$ --- дискретный подграф с 


вершинами $\{ v_1,\dotsc , v_k\}$, то $(T', T^*)$ --- это разлагающий 


фрагмент дерева $T$. Это сразу следует из определений.





Обратно, пусть $(T', T^*)$ --- разлагающий фрагмент $T$. 


Пусть $v_1,\dotsc , v_k$ --- вершины $T^*$. Должна существовать 


вершина $v_0\in V(T')$, $v_0\notin V(T^*)$, cоединенная ребром с одной 


из вершин $T^*$ (а значит, по определению и со всеми). Иначе, т.\,к. 


все вершины $T$, cоединенные ребрами с вершинами из $T^*$, содержатся 


в $T'$, нарушилось бы условие связности $T$. Итак, $v_0$ cуществует, 


но тогда не существует другой такой вершины 


$u_0\in V(T')$, $u_0\notin V(T^*)$, cоединенной ребрами со всеми 


$v_1,\dotsc , v_k$. В противном случае в графе $T$ были бы циклы 


(здесь важно то, что $k\geq 2$). По той же причине не существует 


ребер, соединяющих различные $v_i$ и $v_j$ при $i\neq j, i, j\geq 1$. 


Итак, подграф $T'$ является гроздью, и $T^*$ --- множество его внешних вершин.


\end{exam}





Авторы выражают благодарность рецензенту за полезные замечания.








\begin{thebibliography}{9}   \label{lit}





\bibitem{TK} Тронин С.Н., Копп О.А. {\em Матричные линейные операды} //


Изв. вузов. Математика. -- 2000. -- \No\,6. -- С.\,53--62.





\bibitem{T} Тронин С.Н. {\em Операды конечных графов и гиперграфов} //


Тр. матем. центра им. Н.И.\,Ло\-ба\-чев\-ского. Актуальные проблемы


математики и механики. -- Казань: Унипресс, 2000. -- Т.\,5. -- С.\,207--208.





\bibitem{LM} Levit V.E., Mandrescu E. {\em On hereditary properties


of composition graphs} // Discuss. Math. Graph Theory. -- 1998. -


V.\,18. -- \No\,2. -- P.\,183--195.





\bibitem{H} Харари Ф. {\em Теория графов}. --


М.: Мир, 1973. -- 304~с.





\bibitem{Oz} Ozievich Z. {\em Operad of graphs, convolution


and quasi-Hopf algebra} // Contemp. Math. -- 2003. -- V.\,318.


-- P.\,1--22.








\end{thebibliography}





\vskip 2cm





\post{\it Казанский государственный}{\it Поступила}


\post{\it университет}{28.05.2002}





\label{end}


\end{document}





 








%\bibitem{HP} Харари Ф., Палмер Э. Перечисление графов.---


%М.: "Мир", 1977. -- 328~с.








%\bibitem{G} Гретцер Г. Общая теория решеток.---


%М.: "Мир", 1981. -- 456~с.





%\bibitem{L} Емеличев В.А., Мельников О.И., Сарванов В.И., Тышкевич Р.И.


%Лекции по теории графов.--- М.: "Наука", 1990. -- 384~с.

















%\bibitem{BF} Бордман Дж., Фогт Р. Гомотопически инвариантные


%алгебраические структуры на топологических пространствах . ---


%М.: "Мир", 1977. -- 408~с.











%Покажем на примере, как устроена композиция в операде $\J$.


%Пусть даны графы


%


%\begin{picture}(50,50)


%\put(0,40){\begin{picture}(40,40)%


%\put(10,20){\line(-2,1){10}}


%\put(12,18){3}


%\put(0,25){\line(2,1){14}}


%\put(0,20){1}


%\put(10,20){\line(1,3){4}}


%\put(13,33){2}


%\put(7,9){$\G_1$}


%\end{picture}%}


%\put(20,0){\begin{picture}(20,20)%


%            \put(15,20){\line(-2,3){7}}


%            \put(16,19){2}


%            \put(8,32){1}


%            \put(10,9){$\G_2$}


%         \end{picture}}


%\put(40,0){\begin{picture}(20,20)%


%           \put(15,20){\line(-2,3){7}}


%            \put(16,19){3}


%            \put(6,32){1}


%            \put(8,31){\line(3,1){12}}


%            \put(15,20){\line(1,3){5}}


%            \put(21,34){2}


%            \put(12,9){$\G_3$}


%         \end{picture}}


%\put(70,0){\begin{picture}(20,20)%


%           \put(15,20){\line(-2,3){7}}


%            \put(16,19){2}


%            \put(8,32){1}


%            \put(10,9){$\G_4$}


%         \end{picture}}


%\put(100,0){\begin{picture}(20,20)%


%           \put(15,20){\line(0,1){12}}


%            \put(14,15){3}


%            \put(14,34){1}


%            \put(15,20){\line(1,0){12}}


%            \put(26,15){4}


%            \put(15,32){\line(1,0){12}}


%            \put(26,34){2}


%            \put(17,8){$\G$}


%         \end{picture}}


%


%


%\end{picture}


%


%


%Тогда композицию графов $\G_1\G_2\G_3\G_4 \G$ можно изобразить так (после соответствующей перенумерации):


%


%\begin{picture}(50,50)


%\put(58,45){\line(-2,1){10}}


%\put(58,41){3}


%\put(48,50){\line(2,1){14}}


%\put(47,52){1}


 %\put(30,49){\line(6,1){51}}


 %\put(90,45){\line(-4,1){46}}


 %\put(44,57){\line(1,0){38}}





%\put(58,45){\line(1,3){4}}


%\put(61,58){2}


 %\put(40,45){\line(3,1){45}}


 %\put(40,45){\line(1,0){50}}





%\put(70,0){\begin{picture}(50,50)%


           % \put(20,45){\line(-2,3){8}}


            %\put(21,44){2}


            %\put(13,55){1}


               %\put(40,45){\line(-3,1){8}}


           %\end{picture}}


%\put(25,-10){\begin{picture}(40,40)%


            %\put(15,20){\line(-2,1){15}}


            %\put(16,19){2}


            %\put(-3,29){1}


            


           %\put(15,20){\line(1,5){8}}


           %\put(0,28){\line(6,5){33}} 


           %\put(0,28){\line(3,4){40}}


           %\put(0,28){\line(3,1){40}}





           %\put(15,20){\line(1,2){18}}


           %\put(15,20){\line(3,5){30}}


   


               %\put(15,20){\line(-2,3){10}}


               %\put(16,19){3}


                %\put(6,32){1}


                %\put(8,31){\line(3,1){12}}


                 %\put(15,20){\line(1,3){5}}


                  %\put(21,34){2}


           


         %\end{picture}}


%\put(75,-20){\begin{picture}(40,40)%


           %\put(15,20){\line(-2,3){7}}


            %\put(16,19){3}


            %\put(6,32){1}


            %\put(8,31){\line(3,1){12}}


            %\put(15,20){\line(1,3){5}}


            %\put(21,34){2}








               %\put(15,20){\line(-2,1){20}}


                  %\put(16,19){2}


                     %\put(8,32){1}


           


         %\end{picture}}





%\end{picture}





\





%\begin{figure}[htbp]\label{gr1}


%\setlength{\unitlength}{1 cm}


%\begin{picture}(15,10)(-2.0,-3.0)


%\put(10,10){\special{em:graph Графы.bmp}}


%\end{picture}


%\end{figure}








% Однозначности разложения нет, то есть один и тот же граф допускает


%несколько существенно различных разложений в композицию графов.


%Например:


